\documentclass[11pt]{article}

\usepackage{amsmath,amssymb,amsthm,mathtools}
\usepackage[margin=1in]{geometry}
\usepackage{microtype}
\usepackage[hidelinks]{hyperref}

\newtheorem{theorem}{Theorem}[section]
\newtheorem{lemma}[theorem]{Lemma}
\newtheorem{proposition}[theorem]{Proposition}
\newtheorem{corollary}[theorem]{Corollary}
\newtheorem{remark}[theorem]{Remark}

\newcommand{\dd}{\,\mathrm{d}}
\newcommand{\one}{\mathbf{1}}
\newcommand{\ip}[2]{\left\langle #1,#2\right\rangle}

\title{Cliques with Common Pendant Leaves:\\
A Positive Family for the Chromatic-Polynomial Graphon Bound}
\author{}
\date{}

\begin{document}

\maketitle

\begin{abstract}
Let $H_{r,k}$ be obtained from $K_r$ by attaching $k$ leaves to one
fixed clique vertex.  For every $r\geq3$, $k\geq0$, and every graphon
$W$ of edge density
\[
p\geq1-\frac1{r-1},
\]
we prove
\[
t(H_{r,k},W)
\geq
p^k\prod_{j=0}^{r-1}\bigl(1-j(1-p)\bigr)
=\Phi_{H_{r,k}}(p).
\]
This is the complete density interval required by
$\chi(H_{r,k})=r$.  The proof is a direct link-graphon argument.  It
does not assume a general universal-vertex or join closure.  In the
catalogue on at most six vertices it proves the previously open Atlas
graphs 45, 133, and 191.
\end{abstract}


%%%%%%%%%%%%%%%%%%%%%%%%%%%%%%%%%%%%%%%%%%%%%%%%%%%%%%%%%%%%%%%%%%%%%%%%%%%%%%%
\section{The family and its exact target}
%%%%%%%%%%%%%%%%%%%%%%%%%%%%%%%%%%%%%%%%%%%%%%%%%%%%%%%%%%%%%%%%%%%%%%%%%%%%%%%

Let $(\Omega,\mu)$ be a probability space and let
$W\colon\Omega^2\to[0,1]$ be a symmetric measurable graphon.  For a
finite simple graph $F$, write
\[
t(F,W)
=
\int_{\Omega^{V(F)}}
\prod_{uv\in E(F)}W(x_u,x_v)
\prod_{u\in V(F)}\dd\mu(x_u).
\]
This is the ordinary non-injective homomorphism density.  Put
\[
p=t(K_2,W),
\qquad
d(x)=\int_\Omega W(x,y)\dd\mu(y).
\]

For integers $r\geq3$ and $k\geq0$, define $H_{r,k}$ as follows.
Start with a clique $K_r$, distinguish one of its vertices as $o$, add
new vertices $u_1,\ldots,u_k$, and add precisely the new edges $ou_i$.
Thus all $k$ new vertices are leaves with the same neighbour $o$.

After a proper colouring of $K_r$ has been chosen, each leaf has
$x-1$ choices.  Therefore
\begin{equation}
\chi_{H_{r,k}}(x)=(x)_r(x-1)^k,
\qquad
\chi(H_{r,k})=r,
\label{eq:chromatic}
\end{equation}
where $(x)_r=x(x-1)\cdots(x-r+1)$.  Define
\[
A_r(p):=\prod_{j=0}^{r-1}\bigl(1-j(1-p)\bigr).
\]
If $q=1-p$, then
\begin{align}
\Phi_{H_{r,k}}(p)
&=
q^{r+k}
\left(\frac1q\right)_r
\left(\frac pq\right)^k
\notag\\
&=p^k A_r(p).
\label{eq:target}
\end{align}
The identity extends continuously to $p=1$.


%%%%%%%%%%%%%%%%%%%%%%%%%%%%%%%%%%%%%%%%%%%%%%%%%%%%%%%%%%%%%%%%%%%%%%%%%%%%%%%
\section{The accepted clique input and its tangent lines}
%%%%%%%%%%%%%%%%%%%%%%%%%%%%%%%%%%%%%%%%%%%%%%%%%%%%%%%%%%%%%%%%%%%%%%%%%%%%%%%

We use one previously accepted result: the clique case of the
pure-chordal theorem.  In the notation above it says the following.

\begin{theorem}[Clique-density input]
\label{thm:clique-input}
Let $m\geq2$.  Every graphon $V$ of edge density $z$ satisfying
\[
z\geq a_m:=1-\frac1{m-1}
\]
satisfies
\[
t(K_m,V)\geq A_m(z).
\]
For $m=2$, the convention $a_2=0$ is used, and the assertion is the
identity $t(K_2,V)=z=A_2(z)$.
\end{theorem}

For completeness, the input for $m\geq3$ is exactly the specialization
of the pure-chordal formula to the single maximal clique $K_m$.  It can
also be obtained by multiplying the Moon--Moser clique-ratio inequalities
\[
\frac{t(K_j,V)}{t(K_{j-1},V)}
\geq1-(j-1)(1-z),
\qquad 2\leq j\leq m,
\]
whose right-hand sides are nonnegative when $z\geq a_m$.

\begin{lemma}[Clique-target shape]
\label{lem:shape}
For $m\geq2$, the polynomial
\[
\phi_m(z):=A_m(z)
\]
is nonnegative, nondecreasing, and convex on $[a_m,1]$, and
$\phi_m(a_m)=0$.
\end{lemma}

\begin{proof}
Ignoring the constant factor corresponding to $j=0$, write
\[
\phi_m(z)=\prod_{j=1}^{m-1}\ell_j(z),
\qquad
\ell_j(z)=1-j(1-z).
\]
Every $\ell_j$ is affine and nondecreasing, and every $\ell_j$ is
nonnegative on $[a_m,1]$.  The product is therefore nonnegative.  Its
first derivative is a sum of products of nonnegative factors, and its
second derivative is a sum over ordered pairs of distinct differentiated
factors, again with only nonnegative factors.  Hence
$\phi_m'\geq0$ and $\phi_m''\geq0$ on this interval.  At $a_m$, the
last factor $\ell_{m-1}$ is zero.
\end{proof}

\begin{lemma}[Global tangent bound]
\label{lem:tangent}
Let $m\geq2$ and fix $c\in[a_m,1]$.  Every graphon $V$, of arbitrary
edge density $z\in[0,1]$, satisfies
\[
t(K_m,V)
\geq
\phi_m(c)+\phi_m'(c)(z-c).
\]
\end{lemma}

\begin{proof}
If $z\geq a_m$, Theorem~\ref{thm:clique-input} and convexity give
\[
t(K_m,V)\geq\phi_m(z)
\geq\phi_m(c)+\phi_m'(c)(z-c).
\]
If $z<a_m$, convexity and $\phi_m(a_m)=0$ show that the tangent line
at $c$ is at most zero at $a_m$.  Its slope is nonnegative, so it is
also at most zero at $z<a_m$.  The desired conclusion then follows
from $t(K_m,V)\geq0$.
\end{proof}


%%%%%%%%%%%%%%%%%%%%%%%%%%%%%%%%%%%%%%%%%%%%%%%%%%%%%%%%%%%%%%%%%%%%%%%%%%%%%%%
\section{Weighted link density}
%%%%%%%%%%%%%%%%%%%%%%%%%%%%%%%%%%%%%%%%%%%%%%%%%%%%%%%%%%%%%%%%%%%%%%%%%%%%%%%

Define the rooted triangle density
\[
\tau(x)=
\int_{\Omega^2}
W(x,y)W(x,z)W(y,z)
\dd\mu(y)\dd\mu(z).
\]

\begin{lemma}[Weighted rooted-triangle inequality]
\label{lem:weighted-triangle}
Suppose $p>0$.  For every integer $h\geq0$,
\[
\int_\Omega d(x)^h\tau(x)\dd\mu(x)
\geq
\frac{2p-1}{p}
\int_\Omega d(x)^{h+2}\dd\mu(x).
\]
\end{lemma}

\begin{proof}
Put $M_j=\int d^j$, let $T_W$ be the graphon operator, and write
\[
I_h=\int d^h\tau,
\qquad
J_h=\ip{d^h}{T_Wd}.
\]
The elementary inequality $ab\geq a+b-1$ for $a,b\in[0,1]$, applied
to $W(x,y)$ and $W(x,z)$ and then multiplied by
$d(x)^hW(y,z)$, gives
\begin{equation}
I_h\geq2J_h-pM_h.
\label{eq:I-J}
\end{equation}

Set $U=1-W$, $u=T_U\one=1-d$, and $q=1-p$.  Since
\[
T_Wd=p\one-T_Ud=d-q\one+T_Uu,
\]
we have
\begin{equation}
J_h=M_{h+1}-qM_h+\ip{d^h}{T_Uu}.
\label{eq:J}
\end{equation}
Symmetry of $U$ gives
\begin{align*}
&2\ip{d^h}{T_Uu}
-2\int_\Omega d(x)^hu(x)^2\dd\mu(x)
\\
&\qquad=
\int_{\Omega^2}U(x,y)
\bigl(u(y)-u(x)\bigr)
\bigl(d(x)^h-d(y)^h\bigr)
\dd\mu(x)\dd\mu(y)
\geq0.
\end{align*}
Indeed, $d=1-u$, so $u$ and $d^h$ are oppositely ordered.  Substitution
in \eqref{eq:J}, using $u=1-d$, yields
\[
J_h\geq M_{h+2}-M_{h+1}+pM_h.
\]
Together with \eqref{eq:I-J}, this gives
\[
I_h\geq2M_{h+2}-2M_{h+1}+pM_h.
\]
Finally,
\begin{align*}
&2M_{h+2}-2M_{h+1}+pM_h
-\frac{2p-1}{p}M_{h+2}
\\
&\qquad=
\frac1p\int_\Omega d(x)^h(d(x)-p)^2\dd\mu(x)
\geq0.
\end{align*}
\end{proof}

For $d(x)>0$, define the probability measure
\[
\dd\mu_x(y)=\frac{W(x,y)}{d(x)}\dd\mu(y)
\]
and let $W_x$ denote the same kernel $W$, now viewed on
$(\Omega,\mu_x)$.  Its edge density is
\[
z_x=t(K_2,W_x)=\frac{\tau(x)}{d(x)^2}.
\]
At points where $d(x)=0$, put $z_x=0$.

\begin{corollary}[Weighted average link density]
\label{cor:link-average}
For every integer $n\geq2$ and $p>0$,
\[
\int_\Omega d(x)^nz_x\dd\mu(x)
\geq
\left(2-\frac1p\right)
\int_\Omega d(x)^n\dd\mu(x).
\]
\end{corollary}

\begin{proof}
The left-hand side is $\int d^{n-2}\tau$.  Apply
Lemma~\ref{lem:weighted-triangle} with $h=n-2$.
\end{proof}


%%%%%%%%%%%%%%%%%%%%%%%%%%%%%%%%%%%%%%%%%%%%%%%%%%%%%%%%%%%%%%%%%%%%%%%%%%%%%%%
\section{Main theorem}
%%%%%%%%%%%%%%%%%%%%%%%%%%%%%%%%%%%%%%%%%%%%%%%%%%%%%%%%%%%%%%%%%%%%%%%%%%%%%%%

\begin{theorem}[Cliques with common pendant leaves]
\label{thm:main}
Let $r\geq3$ and $k\geq0$.  Every graphon $W$ with edge density
\[
p\geq1-\frac1{r-1}
\]
satisfies
\[
\boxed{
t(H_{r,k},W)
\geq
p^kA_r(p)
=\Phi_{H_{r,k}}(p).
}
\]
This is the full interval required by $\chi(H_{r,k})=r$.
\end{theorem}

\begin{proof}
Put
\[
m=r-1,
\qquad
n=r-1+k,
\qquad
c=2-\frac1p.
\]
The assumed lower bound on $p$ implies
\[
c\geq1-\frac1{r-2}=a_m.
\]
Also $c\leq1$.  Let $M_n=\int d^n$.

Condition on the image $x$ of the distinguished clique vertex $o$.
The other $r-1$ clique vertices form $K_m$ inside the link graphon
$W_x$, while the $k$ leaves become isolated vertices there.  Directly
from the homomorphism integral,
\begin{equation}
t(H_{r,k},W)
=
\int_\Omega d(x)^n t(K_m,W_x)\dd\mu(x).
\label{eq:link-identity}
\end{equation}
No division is needed on the set $d(x)=0$, because $n\geq2$ and its
contribution to \eqref{eq:link-identity} is zero.

Apply Lemma~\ref{lem:tangent} to each $W_x$, with tangent point $c$:
\[
t(K_m,W_x)
\geq
\phi_m(c)+\phi_m'(c)(z_x-c).
\]
Multiplying by $d(x)^n$ and integrating gives
\begin{align*}
t(H_{r,k},W)
&\geq
M_n\phi_m(c)
\\
&\quad+
\phi_m'(c)
\left(\int d(x)^nz_x\dd\mu(x)-cM_n\right).
\end{align*}
The second line is nonnegative by
Corollary~\ref{cor:link-average} and Lemma~\ref{lem:shape}.  Hence
\[
t(H_{r,k},W)\geq M_n\phi_m(c).
\]
Since $\phi_m(c)\geq0$, Jensen's inequality gives
\[
t(H_{r,k},W)
\geq p^n\phi_m(c).
\]

It remains only to identify this expression.  If $q=1-p$, then
$1-c=q/p$, so
\begin{align*}
\phi_m(c)
&=
\prod_{j=0}^{r-2}\left(1-\frac{jq}{p}\right)
\\
&=
p^{-(r-1)}
\prod_{j=0}^{r-2}\bigl(p-jq\bigr)
\\
&=
p^{-(r-1)}
\prod_{\ell=1}^{r-1}(1-\ell q)
=p^{-(r-1)}A_r(p).
\end{align*}
As $n=r-1+k$, it follows that
\[
p^n\phi_m(c)=p^kA_r(p)=\Phi_{H_{r,k}}(p),
\]
where the last equality is \eqref{eq:target}.
\end{proof}


%%%%%%%%%%%%%%%%%%%%%%%%%%%%%%%%%%%%%%%%%%%%%%%%%%%%%%%%%%%%%%%%%%%%%%%%%%%%%%%
\section{Thresholds and equality}
%%%%%%%%%%%%%%%%%%%%%%%%%%%%%%%%%%%%%%%%%%%%%%%%%%%%%%%%%%%%%%%%%%%%%%%%%%%%%%%

At the required threshold
\[
p=1-\frac1{r-1},
\]
the last factor of $A_r(p)$ is zero.  Thus the target is zero, and the
inequality also follows directly from nonnegativity.  The balanced
complete $(r-1)$-partite graphon attains equality there because it has
no copy of $K_r$.

More generally, for every integer $\ell\geq r-1$, the balanced complete
$\ell$-partite graphon $T_\ell$ has edge density $p=1-1/\ell$ and
\[
t(H_{r,k},T_\ell)
=\frac{\chi_{H_{r,k}}(\ell)}{\ell^{r+k}}
=\Phi_{H_{r,k}}(p).
\]
At $p=1$, one has $W=1$ almost everywhere, and both sides equal $1$.


%%%%%%%%%%%%%%%%%%%%%%%%%%%%%%%%%%%%%%%%%%%%%%%%%%%%%%%%%%%%%%%%%%%%%%%%%%%%%%%
\section{Catalogue consequences}
%%%%%%%%%%%%%%%%%%%%%%%%%%%%%%%%%%%%%%%%%%%%%%%%%%%%%%%%%%%%%%%%%%%%%%%%%%%%%%%

\begin{corollary}[New positive Atlas cases]
The following previously open graphs satisfy the desired inequality on
their full required density intervals:

\begin{center}
\begin{tabular}{ccll}
\hline
Atlas ID & graph6 & family member & target \\
\hline
45  & \texttt{DJ\char123}            & $H_{4,1}$ & $pA_4(p)$ \\
133 & \texttt{E\char126 a?}          & $H_{4,2}$ & $p^2A_4(p)$ \\
191 & \texttt{E\char126\char124?}   & $H_{5,1}$ & $pA_5(p)$ \\
\hline
\end{tabular}
\end{center}

Thus Atlas 45 and 133 are positive for $p\geq2/3$, while Atlas 191 is
positive for $p\geq3/4$.
\end{corollary}

\begin{proof}
Atlas 45 is $K_4$ with one leaf, Atlas 133 is $K_4$ with two leaves
having the same clique neighbour, and Atlas 191 is $K_5$ with one leaf.
The graph6 codes are independent machine-readable identifiers.  Apply
Theorem~\ref{thm:main}.
\end{proof}


%%%%%%%%%%%%%%%%%%%%%%%%%%%%%%%%%%%%%%%%%%%%%%%%%%%%%%%%%%%%%%%%%%%%%%%%%%%%%%%
\section{Scope of the proof}
%%%%%%%%%%%%%%%%%%%%%%%%%%%%%%%%%%%%%%%%%%%%%%%%%%%%%%%%%%%%%%%%%%%%%%%%%%%%%%%

The only previously accepted graph-density result used here is the
clique case of the pure-chordal theorem, stated precisely as
Theorem~\ref{thm:clique-input}.  The weighted rooted-triangle estimate,
the tangent extension, and the link averaging are all proved in this
note.

Although the proof conditions on the distinguished clique vertex, it
does not assert that adjoining a universal vertex preserves positivity
for an arbitrary graph.  It uses the explicit clique polynomial
$A_{r-1}$, its verified convexity, and the exact clique-density input.
Thus it establishes only the stated common-leaf clique family, not a
general join or lifting closure.

All calculations apply directly to arbitrary graphons.  They require
no regularity, finite-step reduction, injective-copy interpretation,
minimum-degree bound, or pointwise lower bound on $d(x)$.

%%%%%%%%%%%%%%%%%%%%%%%%%%%%%%%%%%%%%%%%%%%%%%%%%%%%%%%%%%%%%%%%%%%%%%%%%%%%%%%
\appendix
\section{What the Lean formalization actually proves}
%%%%%%%%%%%%%%%%%%%%%%%%%%%%%%%%%%%%%%%%%%%%%%%%%%%%%%%%%%%%%%%%%%%%%%%%%%%%%%%

Theorem~\ref{thm:main} is formalized as stated, for all $r\geq3$ and $k\geq0$;
Atlas $45$, $133$ and $191$ are \texttt{verified}.  The development is
\texttt{lean/Taeyoung/Methods/CliqueLeaf/}.  The proof follows the note step
for step --- Theorem~\ref{thm:clique-input}, Lemma~\ref{lem:tangent},
Lemma~\ref{lem:weighted-triangle}, Corollary~\ref{cor:link-average} and the
link identity \eqref{eq:link-identity} are each a named theorem there.  Two
remarks.

\subsection*{Theorem~\ref{thm:clique-input} is proved, not accepted}

The clique input is the one previously accepted result the note uses.  It is
formalized, as the clique case of the pure-chordal development
(\texttt{Methods/PureChordal/CliquePolynomialBound.lean}), so no accepted
graph-density result enters the catalogue through this family.

\subsection*{The chain is stated once, abstractly}

Nothing in the argument uses the clique polynomial beyond four facts: it
vanishes at the threshold, it has nonnegative slope, its tangent at $c$ lies
under it on $[a,1]$, and it bounds the base density above the threshold.  The
Lean development therefore states the chain once with the polynomial and the
base graph abstract (\texttt{Methods/ConeBound.lean}), and this family is its
first instance.  Three later families --- the paw and triangle--edge cones, the
verified-base cones, and the clique join of $C_5$ --- are further instances of
the same theorem, so the argument below is not repeated anywhere.

\end{document}
